\problem{P36}{最短随机段与超长段计数}{单纯形平移}
用 $n-1$ 个独立均匀切点把单位线段分成 $n\ge2$ 段，长度为 $D_1,\ldots,D_n$。令 $Z=\min_iD_i$，并对 $t\ge0$ 令 $K_t=\sum_i\ind\{D_i>t\}$。求 $Z$ 的分布、均值和方差，以及 $K_t$ 的概率生成函数与均值。
\solution 切点排序后的联合密度恒为 $(n-1)!$。由切点到间距的线性变换可知，间距向量均匀分布于单纯形 $D_i\ge0,\sum D_i=1$。

若指定 $r$ 个间距均超过 $t$，从这 $r$ 个坐标分别减去 $t$，剩余单纯形的总长度变为 $1-rt$。因其维数为 $n-1$，该事件概率为 $(1-rt)_+^{n-1}$。取 $r=n$，
\[
\boxed{\Prob(Z>z)=(1-nz)^{n-1},\quad0\le z\le1/n.}
\]
支持集之外由 $Z\in[0,1/n]$ 确定。用尾积分及代换 $u=nz$，
\[
\E Z=\frac1{n^2},\qquad
\E Z^2=\frac2{n^3(n+1)},\qquad
\boxed{\Var Z=\frac{n-1}{n^4(n+1)}.}
\]
对计数，展开指示变量乘积：
\[
z^{K_t}=\prod_{i=1}^n[1+(z-1)\ind\{D_i>t\}].
\]
每个 $r$ 元子集有相同的联合概率，故
\[
\boxed{\E z^{K_t}=\sum_{r=0}^n\binom nr(z-1)^r(1-rt)_+^{n-1}.}
\]
此处 $r=0$ 项为 $1$。对 $z$ 在 $1$ 处求导，得到
\[
\boxed{\E K_t=n(1-t)_+^{n-1}.}
\]
\insight 对若干间距同时设下界，等价于平移并缩小同一个单纯形。一次联合概率计算便能同时处理最短段和超阈值段数。
